\section{A misprint in the published family, and its correction}\label{app:typo}

This appendix is not errata history. The point located in \S\ref{sec:point}
lies in a family whose defining matrix is quoted from the published source, and
one entry of that matrix is misprinted. Everything computed in this paper uses
the corrected entry. The correction is therefore part of the mathematics the
paper depends on, and it is set out here in full, with the evidence that
decides it, so that a reader working from the published page reproduces this
work rather than diverging from it silently.

\subsection{The entry, as printed}\label{ssec:typo-printed}

\begin{remark}[The printed reading; \texttt{S16-18}, \textbf{INSPECTED}]
On journal page $2655$, confirmed from the running head, the $b(z)$ entry of
the family prints its first term as
$\bigl(p+\tfrac12(q-1)q\bigr)^{2}$, with \emph{no} $t$, while the second term
carries $(q-1)q\,t$ twice. Read from the same display: $c(z)=-(z-q)/(q-1)$,
hence $c(0)=q/(q-1)$; $a(z)=p+\tfrac12(-2qt+t-\theta_\infty)(z-q)-\tfrac12
t(z-q)^{2}$; ``$q$ is the unique zero of $c(z)$ and $p=a(q)$''; and
``$(p,q)$ gives canonical coordinates on an affine Zariski open set
$U_0=\C\times(\C\setminus\{0,1\})$ of the moduli space $\mathcal M_{t,\lambda}$''.
\end{remark}

The label is \textbf{INSPECTED}: the statement is what the page shows, verified
by looking at a magnified render of it, and it is not the output of a
computation. The same discipline applies to every quotation from the source in
this appendix.

The asymmetry within the entry --- one term with $t$, the neighbouring term
without --- is what first drew attention, but an asymmetry is not an argument.
The argument is that the printed reading fails a test the source itself
supplies.

\subsection{The test that discriminates}\label{ssec:typo-test}

\begin{remark}[The discriminating test, with a control; \texttt{S16-20}, \textbf{VERIFIED}]
Since $A(z)=z(z-1)A_\infty+(z-1)A_0+zA_1$, the source's own constraints
$\det A_i=-\theta_i^{2}/4$ read $\det A(0)=-\theta_0^{2}/4$ and
$\det A(1)=-\theta_1^{2}/4$. The first holds identically for \emph{both}
readings and so does not discriminate. The second holds identically for the
corrected reading $\bigl(p+\tfrac12(q-1)q\,t\bigr)^{2}$, and for the printed
reading leaves the residual
\[
  -\tfrac14 (q-1)(t-1)\bigl(4p+q^{2}t+q^{2}-qt-q\bigr),
\]
which is non-zero at the exact rational point $p=2/7$, $q=5/3$, $t=11/4$,
$\theta_0=\theta_1=0$, $\theta_\infty=3/5$, where it equals $-223/144$. A
deliberately bogus variant $(q-1)q\,t^{2}$ is \emph{rejected} by the same test,
so the test discriminates rather than accepting anything offered to it.
\end{remark}

The control is the part that makes this evidence. A test that accepted the
corrected reading and was silent on everything else would establish only that
the corrected reading is consistent; rejecting a third variant shows the test
has the power to separate.

\begin{remark}[Corroborating structure; \texttt{S16-21}, \textbf{VERIFIED}]
The printed reading's residual carries the factor $(t-1)$ and so vanishes
precisely at $t=1$ --- which is exactly the value at which the two readings
coincide, since $(q-1)q\cdot1=(q-1)q$. A dropped factor of $t$ produces exactly
this signature. \emph{This is corroboration for the diagnosis, not independent
proof of it.}
\end{remark}

The final sentence is the claim's own, and it is reproduced rather than
paraphrased. A signature consistent with a dropped factor is consistent with
other things too; it raises confidence without deciding anything by itself.

\subsection{The same expression, printed correctly, a few lines below}\label{ssec:typo-b1}

The decisive corroboration is on the page itself.

\begin{remark}[The bracket appears twice; \texttt{S16-R18-01}, \textbf{INSPECTED}]
The bracket whose $t$ is missing in $b(z)$ appears a \emph{second} time on
journal page $2655$, in the matrix $B_1$ of the display $B=zB_0+(1/t)B_1$,
where its $(1,2)$ entry prints $\bigl(p+\tfrac12(q-1)q\,t\bigr)^{2}$ \emph{with}
the $t$. This was confirmed by looking at a magnified render. The same
asymmetry is present in \texttt{arXiv:0902.1702v5}, also confirmed by looking.
Both sources therefore print the expression correctly a few lines below the
place where the $t$ is missing.
\end{remark}

\begin{remark}[The two readings, compared exactly; \texttt{S16-R18-02}, \textbf{PROVEN}]
Under the corrected reading, the first term of $b(z)$ equals exactly
$z\cdot B_1[1,2]$; under the printed reading it does not, the difference being
$\tfrac14 z\bigl(4pqt-4pq-4pt+4p+q^{3}t^{2}-q^{3}-2q^{2}t^{2}+2q^{2}+qt^{2}-q\bigr)$.
The difference between the two readings factors as
$-\tfrac14 z(q-1)(t-1)\bigl(4p+q^{2}t+q^{2}-qt-q\bigr)$, carrying the
\emph{same} factor as the $\det A(1)$ residual, so the two independent tests
agree on the correction term for term. Three bogus variants --- $t^{2}$, $t+1$,
$2t$ --- are rejected by the same test.
\end{remark}

Two tests built from different material --- one from the source's determinant
constraints, one from an expression printed twice on the same page --- produce
the same correction factor. That is why this subsection carries the appendix's
only \textbf{PROVEN} label.

\subsection{The misprint persists in the authors' revised version}\label{ssec:typo-arxiv}

\begin{remark}[Present in v5; \texttt{S16-19}, \textbf{INSPECTED}]
\texttt{arXiv:0902.1702} v5, page $34$, prints the identical first term
$\bigl(p+\tfrac12(q-1)q\bigr)^{2}$ with no $t$. The v5 comment records that
``the signs for Hamiltonian systems in Section 4 and the equations (3) (PV) and
(7) (degenerate PV) are corrected'', so that section \emph{was} revised and
this entry was not among the corrections. A search found no published erratum
or corrigendum.
\end{remark}

This matters for reproduction rather than for the mathematics: a reader who
reaches for the most recent version of the source, expecting it to supersede
the journal text, will find the same misprint there.

\subsection{A convention read from the same page}\label{ssec:typo-convention}

One further reading from the same display settles a question that would
otherwise have to be argued structurally.

\begin{remark}[The authors' stated convention; \texttt{S17-17}, \textbf{INSPECTED}]
Page $2655$ states: ``The first basis vector is chosen as cyclic vector and
following Subsection 4.1, $q$ is the unique zero of $c(z)$ and $p=a(q)$.'' This
settles the component question from the authors' own stated convention, and
supersedes an earlier structural argument for the same conclusion while
confirming it. It is also why $q$, being read off $c(z)$, bears on the
second-component reduction rather than the first.
\end{remark}

\subsection{What was adopted, and how}\label{ssec:typo-adopted}

\begin{remark}[The correction, adopted openly; \texttt{S16-05}, \textbf{VERIFIED}]
The corrected reading $\bigl(p+\tfrac12(q-1)q\,t\bigr)^{2}$ is the one used for
computation throughout this work, and it is flagged here rather than adopted
silently. It was discriminated by an exact symbolic residual, not by
preference.
\end{remark}

A closing scope note. Every quotation in this appendix was transcribed from an
image of the page, and no formula here has been re-derived from a computation
that produced it, because there is no such computation --- the source of these
expressions is the printed page. A reader checking this appendix should check
it against the page.
